\newpage
\section{Escritura de Datos}

Las sentencias de escritura forman parte del sublenguaje \textbf{DML}
(\textit{Data Manipulation Language}) de SQL. Permiten insertar, modificar y
eliminar filas en las tablas de una base de datos relacional.

\begin{nota}
  Toda operación DML es transaccional. En motores como MySQL/InnoDB o PostgreSQL,
  los cambios pueden deshacerse con \texttt{ROLLBACK} mientras la transacción
  no haya sido confirmada con \texttt{COMMIT}.
\end{nota}

% ─────────────────────────────────────────────────────────
\subsection{\texttt{INSERT}}

\begin{definicion}{INSERT}
  Agrega \textbf{una o más filas} nuevas a una tabla. Los valores proporcionados
  deben respetar los tipos de datos, restricciones \texttt{NOT NULL} y claves
  foráneas definidas en el esquema.
\end{definicion}

\textbf{Sintaxis — una sola fila:}

\begin{lstlisting}[style=sql, caption={INSERT de una fila con columnas explícitas}]
INSERT INTO nombre_tabla (col1, col2, col3)
VALUES (val1, val2, val3);
\end{lstlisting}

\begin{lstlisting}[style=sql, caption={INSERT sin lista de columnas (no recomendado)}]
-- Los valores deben seguir exactamente el orden de la tabla
INSERT INTO nombre_tabla
VALUES (val1, val2, val3);
\end{lstlisting}

\begin{nota}
  Especificar siempre la lista de columnas hace el código más robusto ante futuros
  cambios en el esquema (columnas nuevas, reordenamiento, etc.).
\end{nota}

\textbf{Sintaxis — múltiples filas:}

\begin{lstlisting}[style=sql, caption={INSERT de varias filas en una sola sentencia}]
INSERT INTO nombre_tabla (col1, col2, col3)
VALUES
    (val1a, val2a, val3a),
    (val1b, val2b, val3b),
    (val1c, val2c, val3c);
\end{lstlisting}

\begin{ejemplo}
  Poblar la tabla \texttt{productos} con tres artículos en una sola operación:

\begin{lstlisting}[style=sql, caption={Inserción múltiple en la tabla productos}]
INSERT INTO productos (nombre, precio, stock, categoria_id)
VALUES
    ('Teclado mecánico',  850.00, 40, 3),
    ('Monitor 24"',      3200.00, 15, 3),
    ('Silla ergonómica', 5400.00,  8, 5);
\end{lstlisting}

  Los campos marcados con \texttt{AUTO\_INCREMENT} (\texttt{producto\_id}) se
  omiten; el motor los genera automáticamente.
\end{ejemplo}

\textbf{Inserción a partir de un \texttt{SELECT}:}

\begin{lstlisting}[style=sql, caption={INSERT ... SELECT: copiar filas desde otra tabla}]
INSERT INTO productos_archivo (nombre, precio, categoria_id)
SELECT nombre, precio, categoria_id
FROM   productos
WHERE  stock = 0;
\end{lstlisting}

% ─────────────────────────────────────────────────────────
\subsection{\texttt{UPDATE}}

\begin{definicion}{UPDATE}
  Modifica los valores de \textbf{una o más columnas} en las filas que satisfacen
  la condición \texttt{WHERE}. Si se omite \texttt{WHERE}, la actualización afecta
  a \textbf{todas} las filas de la tabla.
\end{definicion}

\textbf{Sintaxis general:}

\begin{lstlisting}[style=sql, caption={Sintaxis de UPDATE con WHERE}]
UPDATE nombre_tabla
SET    col1 = nuevo_val1,
       col2 = nuevo_val2
WHERE  condicion;
\end{lstlisting}

\begin{ejemplo}
  Aplicar un incremento de 10\,\% al precio de todos los productos de la
  categoría 3:

\begin{lstlisting}[style=sql, caption={UPDATE con expresión aritmética}]
UPDATE productos
SET    precio = precio * 1.10
WHERE  categoria_id = 3;
\end{lstlisting}
\end{ejemplo}

\begin{ejemplo}
  Actualizar varios campos de un registro específico:

\begin{lstlisting}[style=sql, caption={UPDATE de múltiples columnas por clave primaria}]
UPDATE productos
SET    nombre  = 'Teclado mecánico RGB',
       stock   = 55,
       precio  = 920.00
WHERE  producto_id = 1;
\end{lstlisting}
\end{ejemplo}

\begin{nota}
  \textbf{Buenas prácticas antes de ejecutar un UPDATE:}
  \begin{enumerate}
    \item Verificar la condición con un \texttt{SELECT} previo usando el mismo
          \texttt{WHERE}.
    \item Ejecutar dentro de una transacción explícita para poder hacer
          \texttt{ROLLBACK} si el resultado no es el esperado.
    \item En MySQL, activar \texttt{sql\_safe\_updates = 1} para bloquear
          \texttt{UPDATE} sin \texttt{WHERE} o sin clave en la condición.
  \end{enumerate}
\end{nota}

\textbf{UPDATE con JOIN} (MySQL / MariaDB):

\begin{lstlisting}[style=sql, caption={UPDATE que referencia otra tabla mediante JOIN}]
UPDATE productos AS p
INNER JOIN categorias AS c
    ON p.categoria_id = c.categoria_id
SET p.precio = p.precio * 0.90
WHERE c.nombre = 'Electrónica';
\end{lstlisting}

% ─────────────────────────────────────────────────────────
\subsection{\texttt{DELETE}}

\begin{definicion}{DELETE}
  Elimina las filas que cumplen la condición \texttt{WHERE}. Al igual que
  \texttt{UPDATE}, si se omite \texttt{WHERE} se eliminan \textbf{todas} las
  filas de la tabla, aunque la estructura y los metadatos se conservan.
\end{definicion}

\textbf{Sintaxis general:}

\begin{lstlisting}[style=sql, caption={Sintaxis de DELETE con WHERE}]
DELETE FROM nombre_tabla
WHERE condicion;
\end{lstlisting}

\begin{ejemplo}
  Eliminar los productos sin stock que pertenecen a la categoría 5:

\begin{lstlisting}[style=sql, caption={DELETE con condición compuesta}]
DELETE FROM productos
WHERE stock = 0
  AND categoria_id = 5;
\end{lstlisting}
\end{ejemplo}

\textbf{DELETE con JOIN} (MySQL / MariaDB):

\begin{lstlisting}[style=sql, caption={DELETE que filtra mediante una tabla relacionada}]
DELETE p
FROM   productos AS p
INNER JOIN categorias AS c
    ON p.categoria_id = c.categoria_id
WHERE  c.activa = 0;
\end{lstlisting}

% ── DELETE vs TRUNCATE ────────────────────────────────────
\subsubsection*{Diferencias entre \texttt{DELETE} y \texttt{TRUNCATE}}

\begin{table}[H]
  \centering
  \caption{Comparación entre \texttt{DELETE} y \texttt{TRUNCATE}}
  \renewcommand{\arraystretch}{1.3}
  \rowcolors{2}{gray!10}{white}
  \begin{tabularx}{\textwidth}{>{\bfseries}l X X}
    \toprule
    \textbf{Característica}     & \textbf{\texttt{DELETE}}                        & \textbf{\texttt{TRUNCATE}}                   \\
    \midrule
    Sublenguaje SQL             & DML                                             & DDL                                          \\
    Cláusula \texttt{WHERE}     & Sí — elimina filas selectivamente               & No — elimina todas las filas                 \\
    Transaccional               & Sí — admite \texttt{ROLLBACK}                   & Depende del motor; en MySQL no               \\
    Velocidad                   & Más lenta (registro fila a fila)                & Más rápida (libera páginas de datos)         \\
    Reinicio de \texttt{AUTO\_INCREMENT} & No — el contador no se reinicia       & Sí — el contador vuelve a 1                  \\
    Activación de \textit{triggers} & Sí                                         & No en la mayoría de motores                  \\
    Restricciones de FK         & Respeta \texttt{ON DELETE} definido en el FK   & Falla si existen referencias activas         \\
    \bottomrule
  \end{tabularx}
\end{table}

\begin{lstlisting}[style=sql, caption={Diferencia práctica entre DELETE y TRUNCATE}]
-- Elimina solo las filas del año 2020 (reversible)
DELETE FROM ventas
WHERE YEAR(fecha) = 2020;

-- Vacía la tabla completa y reinicia el auto_increment (irreversible en MySQL)
TRUNCATE TABLE ventas;
\end{lstlisting}

\begin{nota}
  Ante la duda entre \texttt{DELETE} y \texttt{TRUNCATE}, preferir
  \texttt{DELETE} dentro de una transacción explícita: ofrece mayor control y
  trazabilidad. Reservar \texttt{TRUNCATE} para entornos de desarrollo o
  procesos de carga masiva donde la velocidad es crítica y no se requiere
  reversibilidad.
\end{nota}

% ── Resumen del módulo ────────────────────────────────────
\subsection*{Resumen del módulo}

\begin{table}[H]
  \centering
  \caption{Sentencias DML de escritura}
  \renewcommand{\arraystretch}{1.3}
  \rowcolors{2}{gray!10}{white}
  \begin{tabularx}{\textwidth}{>{\ttfamily}l l X}
    \toprule
    \textbf{Sentencia}  & \textbf{Acción}          & \textbf{Consideración clave}                          \\
    \midrule
    INSERT              & Agregar filas            & Especificar siempre la lista de columnas              \\
    INSERT \ldots SELECT & Copiar filas            & Útil para archivado o migración de datos              \\
    UPDATE              & Modificar filas          & Nunca omitir \texttt{WHERE}; verificar con SELECT previo \\
    DELETE              & Eliminar filas           & Transaccional; respeta triggers y claves foráneas     \\
    TRUNCATE            & Vaciar tabla completa    & DDL, más rápido, pero no reversible en MySQL          \\
    \bottomrule
  \end{tabularx}
\end{table}


